% ============================================================
%  Retrieval Results — \input{4_retrieval_results} in main paper
%
%  Required in preamble:
%    \usepackage{booktabs}
%    \usepackage{multirow}
%    \usepackage{xcolor}
%    \usepackage{colortbl}
% ============================================================

% sectiongray for \rowcolor — safe to redefine if main preamble already sets it
\definecolor{sectiongray}{gray}{0.95}

\section{Results}
\label{sec:results}

This section is organised around four research questions:

\begin{description}
  \item[\textbf{RQ1}] Does selecting the segmentation strategy at the question-cluster
    level yield retrieval gains over a single global chunking baseline?
  \item[\textbf{RQ2}] How well do task-conditioned indexing gains transfer to unseen
    questions, and how do they interact with retrieval routing strategies?
  \item[\textbf{RQ3}] Do source-aware retrieval strategies — query routing (QRe) and
    uniform multi-source allocation (UMS) — mitigate the size-imbalance problem in
    heterogeneous multi-source retrieval, and are they complementary components
    targeting structurally distinct failure modes?
\end{description}

\subsection{Task-Conditioned Indexing Evaluation}
\label{sec:indexing-results}

We evaluate whether selecting the segmentation strategy per question cluster
improves over the single-strategy baseline (\texttt{10000\_hierarchical}).
Table~\ref{tab:recall-coverage-k3} reports Recall@3 and Coverage@3 for all
fourteen strategies across all four retrieval methods; no single alternative
outperforms the baseline globally, motivating the cluster-level analysis below.
The cluster-based inference pipeline (Section~\ref{sec:tc-inference-results})
realises a Coverage@3 gain of $+$0.075 ($+$19.2\% relative) on a
held-out test set --- recovering 91.5\% of the oracle upper bound of $+$8.2 points.

\begin{table*}[htbp]
\centering
\resizebox{\textwidth}{!}{%
\begin{tabular}{l cc cc cc cc}
\toprule
\multirow{2}{*}{\textbf{Chunking Strategy}}
  & \multicolumn{2}{c}{\textbf{Naive Retriever @3}}
  & \multicolumn{2}{c}{\textbf{QRe @3}}
  & \multicolumn{2}{c}{\textbf{UMS-Retriever @3}}
  & \multicolumn{2}{c}{\textbf{QRe \& UMS-Retriever @3}} \\
\cmidrule(lr){2-3} \cmidrule(lr){4-5} \cmidrule(lr){6-7} \cmidrule(lr){8-9}
 & Recall & Coverage & Recall & Coverage & Recall & Coverage & Recall & Coverage \\
\midrule
\rowcolor[HTML]{F8F8F8}
\textit{Baseline (10000\_hierarchical)} & \textit{0.362} & \textit{0.362} & \textit{0.481} & \textit{0.480} & \textit{0.504} & \textit{0.502} & \textit{0.525} & \textit{0.523} \\
\midrule
fixed\_10000       & \textbf{0.356} & 0.353          & \textbf{0.475} & \underline{0.470} & \textbf{0.503} & \textbf{0.499} & \textbf{0.525} & \textbf{0.520} \\
S2\_atomic         & 0.188 & 0.221 & 0.308 & 0.342 & 0.395 & 0.417 & 0.416 & 0.437 \\
S3\_window3        & 0.207 & 0.237 & 0.335 & 0.365 & 0.410 & 0.428 & 0.431 & 0.448 \\
S4\_president      & 0.277 & 0.286 & 0.389 & 0.397 & 0.430 & 0.437 & 0.451 & 0.458 \\
S8\_vote           & 0.347 & \textbf{0.373} & 0.444 & \textbf{0.471} & 0.449 & \underline{0.473} & 0.469 & \underline{0.494} \\
S9\_topic          & 0.301 & 0.318 & 0.402 & 0.418 & 0.429 & 0.442 & 0.451 & 0.464 \\
S9\_topic\_noVotes & \underline{0.354} & 0.359 & 0.399 & 0.404 & 0.380 & 0.385 & 0.402 & 0.406 \\
S10\_amendment     & 0.348 & \underline{0.372} & \underline{0.445} & 0.469 & \underline{0.450} & 0.472 & \underline{0.472} & 0.493 \\
S11\_hybrid        & 0.196 & 0.235 & 0.328 & 0.366 & 0.411 & 0.431 & 0.432 & 0.451 \\
S11\_hybrid\_noVotes & 0.203 & 0.239 & 0.331 & 0.367 & 0.411 & 0.431 & 0.433 & 0.451 \\
fixed\_500         & 0.194 & 0.189 & 0.303 & 0.298 & 0.384 & 0.377 & 0.405 & 0.397 \\
fixed\_1000        & 0.172 & 0.181 & 0.293 & 0.301 & 0.381 & 0.383 & 0.402 & 0.403 \\
fixed\_2000        & 0.169 & 0.189 & 0.297 & 0.316 & 0.391 & 0.399 & 0.412 & 0.420 \\
\midrule
\multicolumn{9}{l}{\small \textbf{Note:} The \textit{10000\_hierarchical} reference baseline is excluded from ranking.
  \textbf{Bold}: best; \underline{Underline}: second-best per column.} \\
\bottomrule
\end{tabular}%
}
\caption{Recall and Coverage at $k=3$ for all fourteen chunking strategies across retrieval methods.
No single alternative outperforms \texttt{10000\_hierarchical} globally; gains are cluster-specific.}
\label{tab:recall-coverage-k3}
\end{table*}

%------------------------------------------------------------
\subsubsection{Global vs.\ Cluster-Specific Performance (\textbf{RQ1})}
\label{sec:indexing-global}

Taken in aggregate, no alternative segmentation outperforms \texttt{10000\_hierarchical}:
even \texttt{fixed\_10k} matches it within $<$1 Recall point across all retrieval strategies.
A cluster-level analysis reveals substantial heterogeneity.
Oracle strategy selection, assigning each of the 29 HDBSCAN clusters its
empirically best-performing segmentation at $k{=}3$, yields a weighted average Coverage@3 of 0.471 versus 0.389 for the baseline:
a gain of +8.2 points ($+$21\% relative, $n{=}823$ questions).

The win distribution is highly concentrated: \texttt{S10\_amendment} is optimal
for 13 of the 29 clusters, \texttt{S9\_topic\_noVotes} for 3, and four other
strategies (\texttt{S8\_vote}, \texttt{S4\_president}, \texttt{S3\_window3},
\texttt{fixed\_2000}) for 1--2 clusters each;
the global baseline already wins or ties on the remaining 10 clusters,
including the noise cluster ($n{=}72$).
Table~\ref{tab:chunking-wins-summary} summarises per-strategy gains;
the full per-cluster breakdown is provided in
Appendix~\ref{sec:appendix-clustering-eval}
(Table~\ref{tab:cluster-best-chunking}).

\begin{table}[htbp]
\centering
\small
\setlength{\tabcolsep}{4pt}
\begin{tabular}{l r r r}
\toprule
\textbf{Best chunking} & \textbf{Wins} & \textbf{Median $\Delta$} & \textbf{Max $\Delta$} \\
\midrule
\texttt{S10\_amendment}     & 13 & $+$0.083 & $+$0.222 \\
\texttt{S9\_topic\_noVotes} &  3 & $+$0.083 & $+$0.190 \\
\texttt{S4\_president}       &  1 & $+$0.227 & $+$0.227 \\
\texttt{S3\_window3}         &  1 & $+$0.107 & $+$0.107 \\
\midrule
\textit{Baseline / other}\textsuperscript{$\dagger$}    & 11 & $-$0.018 & $+$0.150 \\
\bottomrule
\end{tabular}
\caption{Cluster wins per chunking strategy at $k{=}3$ (Naive Retriever).
Median and max $\Delta$ are computed across each strategy's winning clusters.
\textsuperscript{$\dagger$}\,The 11-cluster ``Baseline / other'' row aggregates clusters where
\texttt{10000\_hierarchical} is already optimal or where the
``best'' alternative shows no positive improvement (\texttt{fixed\_10000},
\texttt{S8\_vote}, \texttt{S9\_topic}, \texttt{fixed\_2000}).
Full per-cluster breakdown: Appendix Table~\ref{tab:cluster-best-chunking}.}
\label{tab:chunking-wins-summary}
\end{table}

%------------------------------------------------------------
\subsubsection{Which Strategies Gain and Why}
\label{sec:indexing-analysis}

\paragraph{\texttt{S10\_amendment} (13 wins, $\Delta$ up to $+$0.222).}
Amendment lifecycle segmentation captures the full argumentative arc of a
legislative proposal — proposer's argument, opposition response, minister's
defence, closing vote — as one retrievable unit.
The baseline merges multiple amendment cycles into a single 10k chunk, diluting
the relevant passage.
Largest gains for this strategy: cluster~15 ($+$0.222), cluster~7 ($+$0.179),
cluster~25 ($+$0.143).
The overall maximum $\Delta$ ($+$0.227, cluster~4; Appendix
Table~\ref{tab:cluster-best-chunking}) belongs to \texttt{S4\_president},
which isolates presidential utterances for clusters centred on
speaker-attribution questions.

\paragraph{\texttt{S9\_topic\_noVotes} (3 wins).}
Clusters requiring synthesis over a full agenda item benefit from topic-level
segmentation.
Removing roll-call vote name-lists (\textit{``ONT VOT\'{E} POUR: MM.~\ldots''})
is critical: these blocks are semantically orthogonal to debate content and
contaminate query-document similarity.
Notably, \texttt{S9\_noVotes} is \emph{harmed} by UMS-Retriever (Recall .380
vs.\ .354 with Naive), suggesting its reduced chunk count interacts poorly with
equal-quota allocation.

\paragraph{Clusters where the baseline wins (10 clusters).}
Questions requiring co-retrieval of a newspaper article and a long parliamentary
extract, sustained ministerial monologues, or topically heterogeneous sessions
all favour the 10k context window.
%Per-cluster qualitative discussion is in Appendix~\ref{sec:appendix-cluster-findings}.

%------------------------------------------------------------
\subsubsection{Task-Conditioned Inference on Held-Out Test Set (\textbf{RQ1, RQ2})}
\label{sec:tc-inference-results}

The cluster-based inference pipeline (Algorithm~\ref{alg:tc-inference}) is
evaluated on the held-out 25\% test split ($n_{\text{test}}{=}221$ questions,
seed\,=\,42), applying the strategy map $s^*$ learned on the training split to
unseen queries via nearest-centroid assignment in UMAP space.
All four retrieval strategies (Naive, QRe, UMS, QRe\,\&\,UMS) are evaluated
under both chunking policies; each TC condition is trained and evaluated under
the same routing function to avoid strategy-selection mismatch.
Table~\ref{tab:tc-inference} reports Coverage@$k$ for $k \in \{3,5,10\}$.

Without routing, task-conditioned chunking yields the largest gains. Under the Naive Retriever, Coverage@3 improves by $+$0.075 ($+$19.2\% relative)
on unseen questions, rising to $+$0.081 at $k{=}5$.
The realised $+$0.075 Coverage@3 represents \textbf{91.5\% of the oracle gap}
of $+$0.082 (Appendix Table~\ref{tab:cluster-best-chunking}: oracle $0.389 \to 0.471$;
realised $0.390 \to 0.465$), confirming that UMAP nearest-centroid assignment
transfers training-phase strategy preferences with minimal loss to unseen queries.

Routing strategies reduce and can reverse the TC gain. Under QRe, the Coverage@3 advantage shrinks to $+$0.016 and vanishes at larger
budgets ($+$0.001 at @5, $-$0.012 at @10).
Under UMS and QRe\,\&\,UMS, task-conditioned chunking produces a small decrease at all cut-offs ($-$0.005 to $-$0.013).
This interaction is expected: routing already removes the most irrelevant
candidates from the pool, narrowing the margin in which chunking granularity
can contribute additional signal.

\begin{table*}[htbp]
\centering
\small
\setlength{\tabcolsep}{4pt}
\begin{tabular}{ll ccc}
\toprule
\textbf{Retrieval} & \textbf{Chunking} & \textbf{Cov@3} & \textbf{Cov@5} & \textbf{Cov@10} \\
\midrule
\multirow{3}{*}{Naive}
  & Global Baseline  & 0.390 & 0.457 & 0.571 \\
  & Task-Conditioned & \textbf{0.465} & \textbf{0.538} & \textbf{0.638} \\
  & $\Delta$         & \cellcolor{db1}{+0.075} & \cellcolor{db1}{+0.081} & \cellcolor{db1}{+0.067} \\
\midrule
\multirow{3}{*}{QRe}
  & Global Baseline  & 0.495 & 0.578 & \textbf{0.687} \\
  & Task-Conditioned & \textbf{0.511} & \textbf{0.579} & 0.675 \\
  & $\Delta$         & \cellcolor{db1}{+0.016} & +0.001 & $-$0.012 \\
\midrule
\multirow{3}{*}{UMS}
  & Global Baseline  & \textbf{0.534} & \textbf{0.598} & \textbf{0.715} \\
  & Task-Conditioned & 0.527 & 0.591 & 0.702 \\
  & $\Delta$         & $-$0.007 & $-$0.007 & $-$0.013 \\
\midrule
\multirow{3}{*}{QRe \& UMS}
  & Global Baseline  & \textbf{0.546} & \textbf{0.607} & \textbf{0.731} \\
  & Task-Conditioned & 0.541 & 0.598 & 0.718 \\
  & $\Delta$         & $-$0.005 & $-$0.009 & $-$0.013 \\
\bottomrule
\end{tabular}
\caption{Coverage@$k$ on the held-out test set ($n_{\text{test}}{=}221$):
Global Baseline (\texttt{10000\_hierarchical}) vs Task-Conditioned chunking,
across all four retrieval strategies.
Each TC condition is trained and evaluated under the same routing function.
Positive $\Delta$ shaded light blue; task-conditioned gains are concentrated
under the Naive Retriever and diminish when routing is active.}
\label{tab:tc-inference}
\end{table*}

% \begin{figure}[t]
% \centering
% \includegraphics[width=\linewidth]{images/tc_conditions_bar.pdf}
% \caption{Coverage@$k$ for $k\in\{3,5,10\}$: Global Baseline (blue) vs
% Task-Conditioned (green) across all four retrieval strategies.
% $\Delta$ annotations show the absolute gain; task-conditioned improvements are
% concentrated under the Naive Retriever and diminish when routing is active.}
% \label{fig:tc-bar}
% \end{figure}

% Cluster-level generalisation analysis (per-cluster training vs test Coverage@3
% diagonal plot) moved to Appendix~\ref{sec:appendix-clustering-eval}.

%------------------------------------------------------------
\subsubsection{Practical Implications (\textbf{RQ2})}
\label{sec:indexing-practical}

The oracle gain of $+$8.2 points establishes an upper bound for task-conditioned
indexing; the nearest-centroid inference of Algorithm~\ref{alg:tc-inference}
realises \textbf{91.5\%} of that upper bound ($+$7.5 Coverage@3) on unseen
questions.
All indexes are pre-built offline; the only marginal online cost is a single
UMAP projection and a nearest-neighbour lookup over $|\{\mu_c\}| = 29$ centroids
— negligible relative to the query-embedding step.


%------------------------------------------------------------
\subsection{Retrieval Evaluation}
\label{sec:retrieval-results}



We evaluate the four retrieval strategies described in
Section~\ref{sec:retrieval-strategies} on the multi-hop subset of
HistoriQA-ThirdRepublic~\cite{historiqa}.
Results are broken down by question type — \textit{cross-newspaper} questions (both
gold documents from newspaper corpora, 314 questions) and
\textit{newspaper $\rightarrow$ Débats} questions (one gold document from a newspaper
and one from parliamentary transcripts, 571 questions) — to expose the asymmetric
effects of each strategy.
HippoRAGv2~\cite{hipporagv2} serves as a knowledge-graph-based external baseline.

%------------------------------------------------------------

\begin{table*}[ht]
\centering
\caption{Retrieval performance (Recall@$k$) on multi-hop questions, broken down by
question type.
\textbf{S-RAG}: Structured RAG pipeline proposed in this work.
\textbf{QRe}: Query Rerouting — supervised source pre-selection.
\textbf{UMS}: Uniform Multi-Source Retriever — equal-quota retrieval across collections.
\textbf{Bold} = best; \underline{underline} = second-best per column.}
\label{tab:retrieval-results}
\small
\setlength{\tabcolsep}{0pt}
\begin{tabular*}{\textwidth}{@{\extracolsep{\fill}} l ccc @{}}
\toprule
\textbf{Configuration}
  & \textbf{Recall@3}
  & \textbf{Recall@5}
  & \textbf{Recall@10} \\
\midrule

\rowcolor{sectiongray}
\textbf{Naive Retriever}~\cite{historiqa} & 35.9 & 44.0 & 54.4 \\
\quad Cross-newspaper (314 q.)            & 14.3 & 19.7 & 28.8 \\
\quad Newspaper $\rightarrow$ Débats (571 q.) & 47.8 & 57.3 & 68.5 \\

\midrule
\rowcolor{sectiongray}
\textbf{HippoRAGv2}~\cite{hipporagv2}       & 31.7 & 41.9 & 53.0 \\
\quad Cross-newspaper (314 q.)            & 11.6 & 17.2 & 24.7 \\
\quad Newspaper $\rightarrow$ Débats (571 q.) & 42.7 & 55.4 & 68.6 \\

\midrule
\rowcolor{sectiongray}
\textbf{S-RAG : QRe}                     & 48.0 & 57.0 & 68.4 \\
\quad Cross-newspaper (314 q.)            & 48.4 & 56.7 & 68.3 \\
\quad Newspaper $\rightarrow$ Débats (571 q.) & 47.8 & 57.2 & 68.5 \\

\rowcolor{sectiongray}
\textbf{S-RAG : UMS-Retriever}           & \underline{50.2} & \underline{57.6} & \underline{69.4} \\
\quad Cross-newspaper (314 q.)            & 40.8 & 53.3 & 60.7 \\
\quad Newspaper $\rightarrow$ Débats (571 q.) & 55.4 & 60.0 & 74.2 \\

\midrule
\rowcolor{sectiongray}
\textbf{S-RAG : QRe \& UMS-Retriever}   & \textbf{52.3} & \textbf{58.4} & \textbf{72.3} \\
\quad Cross-newspaper (314 q.)            & 46.7 & 55.6 & 69.3 \\
\quad Newspaper $\rightarrow$ Débats (571 q.) & 55.4 & 60.0 & 73.9 \\

\bottomrule
\end{tabular*}
\end{table*}

%------------------------------------------------------------
\subsubsection{Overall Gains over Baseline (\textbf{RQ3})}
\label{sec:results-overall}

All three proposed strategies substantially outperform both the Naive Retriever
and HippoRAGv2 across all cut-offs.
The combined QRe \& UMS-Retriever achieves Recall@3 of 52.3, a gain of +16.4 points over the naive baseline (+46\% relative) and +20.6 points over HippoRAGv2.
At Recall@10, the gap widens further (+17.9 pts vs.\ naive), confirming that
the structured strategies succeed in surfacing more gold documents even with
a larger candidate window.

These results establish that neither routing alone nor uniform allocation alone
accounts for the full gain: each component contributes independently, and their
combination is strictly dominant on every metric.

%------------------------------------------------------------
\subsubsection{Asymmetric Effects by Question Type (\textbf{RQ3})}
\label{sec:results-asymmetric}

The per-type breakdown reveals a clear asymmetry in how the two strategies operate.

\paragraph{Cross-newspaper questions.}
The Naive Retriever achieves only 14.3 Recall@3 on these 314 questions, reflecting
the well-known dominance problem: the large \textit{Les Débats} collection floods
the top results for queries that do not require parliamentary evidence at all.
QRe eliminates this noise by routing the query away from \textit{Les Débats},
boosting Recall@3 to 48.4 ($\times$3.4 relative improvement).
By contrast, UMS-Retriever on its own is less effective here (40.8), because
it still allocates one third of the budget to \textit{Les Débats};
QRe's source filtering is the decisive factor for this question type.

\paragraph{Newspaper $\rightarrow$ Débats questions.}
On the 571 questions that require cross-source evidence, the Naive Retriever already
performs well (47.8 Recall@3) due to the naturally high relevance of parliamentary
passages for these queries.
UMS-Retriever provides the main improvement here (55.4 Recall@3, +7.6 pts),
by guaranteeing that the \textit{Les Débats} collection always contributes at
least one third of the retrieved pool rather than being crowded out by
high-scoring newspaper passages.
QRe's benefit on this split is marginal (47.8 unchanged at @3), since all three
collections are already active under the ``otherwise'' branch of the routing map.

%------------------------------------------------------------
\subsubsection{Complementarity of QRe and UMS (\textbf{RQ3})}
\label{sec:results-complementarity}

The two strategies target structurally different failure modes of the Naive Retriever:
QRe corrects \emph{source contamination} (irrelevant collections dominating the pool),
while UMS-Retriever corrects \emph{source starvation} (relevant collections
under-represented despite high relevance).
These failure modes are not mutually exclusive — both can occur simultaneously for
a given question — which explains why combining them yields strictly better results
than either in isolation.
At Recall@5, QRe alone achieves 57.0, UMS alone achieves 57.6, but their combination
reaches 58.4, with consistent improvements across both question types.

\paragraph{Comparison to HippoRAGv2.}
HippoRAGv2, despite leveraging a rich graph-based index, underperforms the Naive
Retriever on cross-newspaper questions (11.6 vs.\ 14.3 @3) and falls short
of all three proposed strategies on every metric.
This suggests that graph-propagation methods may struggle when evidence is distributed
across multiple heterogeneous corpora with different writing styles and registers,
a setting our structured multi-source approach is explicitly designed for.


% \todo{Add Answer Generation Evaluation section --- end-to-end QA results with best retrieval config, ablation, error analysis.}

% \todo{Add error analysis --- qualitative failure examples: cross-source bridging, OCR mismatches, long-doc compression.}